\documentclass[11pt,a4paper]{article}
\usepackage[utf8]{inputenc}
\usepackage[T1]{fontenc}
\usepackage{lmodern}
\usepackage{textcomp}
\usepackage{amsmath,amssymb}
\usepackage{graphicx}
\usepackage{booktabs}
\usepackage{longtable}
\usepackage{array}
\usepackage{listings}
\usepackage{xcolor}
\usepackage[margin=1in]{geometry}
\usepackage{float}
\usepackage[hidelinks,breaklinks]{hyperref}
\usepackage{microtype}
\setlength{\parskip}{4pt}
\lstset{basicstyle=\ttfamily\scriptsize,breaklines=true,breakatwhitespace=false,
  postbreak=\mbox{\textcolor{gray}{$\hookrightarrow$}\space},frame=single,
  framerule=0.3pt,rulecolor=\color{gray!60},columns=fullflexible,keepspaces=true,
  showstringspaces=false,numbers=none,xleftmargin=4pt,xrightmargin=4pt,
  literate={\_}{{\_}}1}
\graphicspath{{figures/}}
\providecommand{\tightlist}{}
\title{Transcript: The Hubble Tension: A Conceptual Summary of Early- and Late-Universe Determinations of $H_0$}
\author{HEP-Genesis-Agent session record \\ \small Project: ramp-features}
\date{2026-10-15}
\begin{document}
\maketitle
\noindent\textit{\small The conversation verbatim, generated mechanically from the session record.}
\vspace{1em}



\appendix
\section{Conversation transcript}
\label{app:transcript}
The exchange verbatim (assistant text as shown to the user; tool calls summarized by name --- see the tool-call log for arguments and results).
\subsection*{User \small(2026-10-15 22:42:59)}
What is the Hubble tension? Two paragraphs, then write a report.

\subsection*{Assistant \small(quick mode, 2026-10-15 22:42:59)}
The Hubble tension is a disagreement between two ways of measuring the universe's current expansion rate, the Hubble constant H0. In the ''early-universe'' method, you fit the standard $\Lambda$CDM model (dark energy plus cold dark matter) to the cosmic microwave background and project forward to today. The Planck satellite gets H0 = 67.4 $\pm$ 0.5 km/s/Mpc this way (Planck Collaboration 2018, arXiv:1807.06209). Baryon acoustic oscillations and Big Bang nucleosynthesis give a similar value. In the ''late-universe'' method, you build a distance ladder: parallaxes and geometric anchors calibrate Cepheid variable stars, the Cepheids calibrate Type Ia supernovae, and the supernovae measure distances deep into the Hubble flow. The SH0ES team gets H0 = 73.04 $\pm$ 1.04 km/s/Mpc this way (Riess et al. 2022, arXiv:2112.04510). The two values differ by about 5$\sigma$. That is too big a gap to blame on chance if the quoted uncertainties are right.

There are two broad ways to explain it, and neither has won. One is a systematic error in the measurements. Candidates include crowding and photometry problems with Cepheids, the metallicity dependence of Cepheid brightness, or how the supernovae are standardized. Calibrating with the tip of the red giant branch instead of Cepheids gives intermediate values near 69.8 km/s/Mpc (Freedman et al. 2019, arXiv:1907.05922), which keeps that debate going. The other is new physics beyond $\Lambda$CDM, such as early dark energy, extra relativistic species, changes to recombination, or dark-sector interactions. Most of these proposals ease the H0 discrepancy only by making fits to other data worse, such as galaxy clustering (S8) or BAO (see the review by Di Valentino et al. 2021, arXiv:2103.01183). Independent methods may settle it: gravitational-wave ''standard sirens'', time delays in strongly lensed quasars, and new JWST observations of Cepheids and red giants.

These citations come from my own knowledge, not from a literature search in this session. The values are as published in the papers named. I'll now write the report.

\end{document}
